\scanpage{383}
\mspage{299}
% archive sheet number: 382

il existe une application
\[
\unsure{\varphi|\phi}:I_{\alpha_\pi}\unsure{\longrightarrow|\text{---}}A
\]
telle que pour tout $r\in A$, $\varphi^{-1}(\{r\})$ soit
une partie cofinale de $I_{\alpha_\pi}$, ou ce
qui \deleted{revient au $=$}\inserted{est plus \deleted{\unsure{fort}}précis},
$\operatorname{card}\varphi^{-1}(\{r\})=\pi$.

En effet il suffit de considérer
un ensemble $A'$ de cardinal $\pi$. (p. ex. $A'=I_{\alpha_\pi}$), et
l'application
\[
\psi=\operatorname{pr}_1:A\times A'\unsure{\text{---}|\longrightarrow}A,
\]
et de se donner une bijection \deleted{de}
\[
\deleted{A\times A}\qquad A\times A'\xrightarrow{\sim}I_{\alpha_\pi},
\]
laquelle existe puisque
\[
\operatorname{card}(A\times A')=\operatorname{card}A\mathbin{.}\operatorname{card}A'=\operatorname{card}A'=\pi\inserted{{}=\operatorname{card}I_{\alpha_\pi}}
\]
puisque $0\neq\operatorname{card}A\leq\operatorname{card}A'=\pi$ et \deleted{\unsure{$\operatorname{card}A'$ infini}}\inserted{$\pi$ infini}.

\deleted{On choisira} On suppose choisie $\varphi$
comme dans 4.13.10.1, et on \deleted{choisi\unsure{ra}}
construira le système inductif
$(i_\alpha)_{\alpha\in\deleted{\unsure{A}}I_{\alpha_\pi}}$ de telle façon qu'on ait
\[
i_\alpha\in J_{\varphi(\alpha)}.\tag*{\deleted{$\unsure{(2)}$} $(*)$}
\]
Posant
\[
j=\varinjlim_{\alpha\in I_{\unsure{\alpha|\alpha_\pi}}}i_\alpha,
\]
